\section{Data Collection and Preparation}
\label{sec:3_2}

\subsection{Data Sources}

The dataset comprises two primary sources from SMG's production systems covering the period from January 2024 to November 2025.

\textbf{Insertion Events} (\texttt{insertion\_events.csv}) contains 804,717 lifecycle events for real estate listings across Homegate and ImmoScout24 platforms. Each event records a status change (DRAFT, PENDING\_APPROVAL, APPROVED, PUBLISHED, ARCHIVED, DELETED) with associated listing attributes including property details, pricing, geographic information, and lister contact data.

\textbf{SEON Transactions} (\texttt{seon\_transactions.csv}) contains approximately 140,000 fraud detection API responses capturing device fingerprints, IP intelligence, email validation, phone verification, and behavioral signals collected at listing submission time.

\textbf{Data Merge Strategy:} Events and SEON transactions occur at different timestamps, requiring temporal alignment. We implement an as-of join strategy:

\begin{enumerate}
    \item For each event at time T, find the closest SEON transaction after time T (preferred).
    \item If no subsequent transaction exists, use the closest SEON transaction before time T.
\end{enumerate}

This approach ensures features reflect information available at decision time, simulating realistic operational conditions where SEON data may arrive slightly before or after the event.

\subsection{Data Cleaning and Preprocessing}

\textbf{Datetime Handling:} The two data sources use different timestamp formats requiring standardization:
\begin{itemize}
    \item Events: \texttt{\%Y-\%m-\%d \%H:\%M:\%S\%.f \%z} (e.g., ``2025-01-15 10:30:45.123 +0100'')
    \item SEON: \texttt{\%Y-\%m-\%dT\%H:\%M:\%S\%.f\%z} (e.g., ``2025-01-15T10:30:45.123+0000'')
\end{itemize}

All timestamps are converted to UTC for consistent temporal operations.

\textbf{PII Anonymization:} Personally identifiable information is hashed using SHA-256 before analysis:
\begin{itemize}
    \item \texttt{INSERTION\_ID} $\rightarrow$ \texttt{INSERTION\_ID\_hash}
    \item \texttt{LISTING\_LISTER\_EMAIL} $\rightarrow$ \texttt{LISTING\_LISTER\_EMAIL\_hash}
    \item \texttt{LISTING\_LISTER\_PHONE} $\rightarrow$ \texttt{LISTING\_LISTER\_PHONE\_hash}
    \item \texttt{user\_id} $\rightarrow$ \texttt{user\_id\_hash}
    \item IP addresses $\rightarrow$ \texttt{ip\_hash}
    \item Device fingerprints $\rightarrow$ \texttt{session/device\_hash}
\end{itemize}

\textbf{Temporal Filtering:} Analysis of data density revealed sparse coverage before November 2024, when SEON integration commenced. Events from January-October 2024 average approximately 200 events per month compared to 60,000+ post-integration. We filter to \texttt{train\_start\_date = 2024-12-01} to ensure consistent feature availability across the training period.

\textbf{ETL Implementation:} The pipeline uses Polars for memory-efficient processing of the 800K+ event dataset, with lazy evaluation enabling streaming operations on commodity hardware.

\subsection{Label Definition}

\textbf{Fraud Label Source:} The \texttt{FLAGGEDFORFRAUD} column contains a timestamp indicating when customer support or automated rules flagged a listing as fraudulent. Label assignment follows:
\begin{itemize}
    \item If \texttt{FLAGGEDFORFRAUD} is non-null $\rightarrow$ \texttt{is\_fraud = 1}
    \item If \texttt{FLAGGEDFORFRAUD} is null $\rightarrow$ \texttt{is\_fraud = 0}
\end{itemize}

\textbf{Critical Observation:} Fraud is a property of listings, not individual events. A single fraudulent listing may generate multiple events (DRAFT $\rightarrow$ PENDING\_APPROVAL $\rightarrow$ DELETED) before detection. Analysis must aggregate to listing level to avoid inflating sample counts.

\subsection{Label Leakage Discovery and Mitigation}

During comprehensive data mining analysis, we discovered a critical label leakage issue in the \texttt{STATUS} column that illustrates the importance of temporal reasoning in fraud detection model development.

\textbf{Discovery Process:} Initial feature importance analysis revealed unexpectedly high predictive power for the \texttt{STATUS} feature. Further investigation of the event lifecycle revealed the following patterns:

\begin{itemize}
    \item \textbf{Fraudulent listings:} DRAFT $\rightarrow$ PENDING\_APPROVAL $\rightarrow$ DELETED (flagged=True)
    \item \textbf{Legitimate listings:} DRAFT $\rightarrow$ PENDING\_APPROVAL $\rightarrow$ APPROVED $\rightarrow$ PUBLISHED $\rightarrow$ ARCHIVED
\end{itemize}

\textbf{Root Cause:} The leakage stems from a temporal misalignment between prediction time and data collection:
\begin{enumerate}
    \item At \textbf{prediction time} (listing submission), all listings have \texttt{STATUS = DRAFT}
    \item In \textbf{our dataset} (historical snapshot), fraudulent listings show \texttt{STATUS = DELETED}
    \item The deletion occurs \textbf{after} fraud detection, not before
    \item A model learning ``DELETED $\rightarrow$ FRAUD'' exploits backwards causality
\end{enumerate}

\textbf{Impact Analysis:} Initial experiments including STATUS as a feature showed artificially inflated performance metrics. The model was effectively learning post-hoc outcomes rather than predictive patterns available at decision time. This would result in catastrophic deployment failure, as the STATUS signal would not be available for real-time predictions.

\textbf{Mitigation:} Following this discovery, the \texttt{STATUS} column was permanently removed from the feature schema defined in \texttt{src/features/schema.py}. All reported results in this thesis reflect models trained without STATUS, ensuring fair evaluation and realistic deployment expectations.

This discovery reinforces the critical importance of point-in-time feature engineering in temporal prediction tasks, where the distinction between information available at prediction time versus information available in retrospective analysis must be rigorously maintained.
